\documentclass{article}

% packages
\usepackage{amsmath, amsthm, thmtools, amsfonts, amssymb, luacode, catchfile, tikzducks, hyperref, ifthen}
\ifcsname c@kobocompile\endcsname
	\usepackage[a5paper, total={1072pt, 1448pt}, margin=10pt, includeheadfoot]{geometry} % set page margins
\else
	\usepackage[a4paper, margin=50pt, includeheadfoot]{geometry}
\fi
\usepackage[shortlabels]{enumitem}
\usepackage[skip=3pt, indent=0pt]{parskip}

% language
\usepackage[bidi=basic, layout=tabular, provide=*]{babel}
\ifcsname c@english\endcsname
	\babelprovide[main, import]{english}
\else
	\babelprovide[main, import]{hebrew}
	\babelprovide{rl}
\fi
%\babelfont{rm}{Libertinus Serif}
\babelfont{rm}[Renderer=Harfbuzz]{Libertinus Serif}
\babelfont{sf}{Libertinus Sans}
\babelfont{tt}{Libertinus Mono}

% style
\AddToHook{cmd/section/before}{\clearpage}	% Add line break before section
\linespread{1.3}
\setcounter{secnumdepth}{0}		% Remove default number tags from sections, this won't do well with theorems
\AtBeginDocument{\setlength{\belowdisplayskip}{3pt}}
\AtBeginDocument{\setlength{\abovedisplayskip}{3pt}}
\graphicspath{ {../images/} }

% operators
\DeclareMathOperator\cis{cis}
\DeclareMathOperator\Sp{Sp}
\DeclareMathOperator\tr{tr}
\DeclareMathOperator\im{Im}
\DeclareMathOperator\re{Re}
\DeclareMathOperator\diag{diag}
\DeclareMathOperator*\lowlim{\underline{lim}}
\DeclareMathOperator*\uplim{\overline{lim}}
\DeclareMathOperator\rng{rng}
\DeclareMathOperator\Sym{Sym}
\DeclareMathOperator\Arg{Arg}
\DeclareMathOperator\Log{Log}
\DeclareMathOperator\dom{dom}
\DeclareMathOperator\supp{Supp}
\DeclareMathOperator\var{Var}
\DeclareMathOperator\cov{Cov}

% commands
%\renewcommand\qedsymbol{\textbf{מש''ל}}
%\renewcommand\qedsymbol{\fbox{\emoji{lizard}}}
\newcommand{\Aa}[0]{\mathcal{A}}
\newcommand{\Bb}[0]{\mathcal{B}}
\newcommand{\CC}[0]{\mathbb{C}}
\newcommand{\Cc}[0]{\mathcal{C}}
\newcommand{\EE}[0]{\mathbb{E}}
\newcommand{\FF}[0]{\mathbb{F}}
\newcommand{\Ff}[0]{\mathcal{F}}
\newcommand{\Ii}[0]{\mathcal{I}}
\newcommand{\Gg}[0]{\mathcal{G}}
\newcommand{\Ll}[0]{\mathcal{L}}
\newcommand{\Mm}[0]{\mathcal{M}}
\newcommand{\NN}[0]{\mathbb{N}}
\newcommand{\Nn}[0]{\mathcal{N}}
\newcommand{\PP}[0]{\mathbb{P}}
\newcommand{\Pp}[0]{\mathcal{P}}
\newcommand{\QQ}[0]{\mathbb{Q}}
\newcommand{\RR}[0]{\mathbb{R}}
\newcommand{\Rr}[0]{\mathcal{R}}
\newcommand{\Ss}[0]{\mathcal{S}}
\newcommand{\TT}[0]{\mathbb{T}}
\newcommand{\Uu}[0]{\mathcal{U}}
\newcommand{\Vv}[0]{\mathcal{V}}
\newcommand{\Ww}[0]{\mathcal{W}}
\newcommand{\ZZ}[0]{\mathbb{Z}}
\newcommand{\acts}[0]{\circlearrowright}
\newcommand{\explain}[2] {
	\begin{flalign*}
		 && \text{#2} && \text{#1}
	\end{flalign*}
}
\newcommand{\maketitleprint}[0]{ \begin{center}
	%\begin{tikzpicture}[scale=3]
	%	\duck[graduate=gray!20!black, tassel=red!70!black]
	%\end{tikzpicture}	
	\includegraphics[width=6cm]{cover}
\end{center}
}

% theorem commands
\newtheoremstyle{c_remark}
	{}	% Space above
	{}	% Space below
	{}% Body font
	{}	% Indent amount
	{\bfseries}	% Theorem head font
	{}	% Punctuation after theorem head
	{.5em}	% Space after theorem head
	{\thmname{#1}\thmnumber{ #2}\thmnote{ \normalfont{\text{(#3)}}}}	% head content
\newtheoremstyle{c_definition}
	{3pt}	% Space above
	{3pt}	% Space below
	{}% Body font
	{}	% Indent amount
	{\bfseries}	% Theorem head font
	{}	% Punctuation after theorem head
	{.5em}	% Space after theorem head
	{\thmname{#1}\thmnumber{ #2}\thmnote{ \normalfont{\text{(#3)}}}}	% head content
\newtheoremstyle{c_plain}
	{3pt}	% Space above
	{3pt}	% Space below
	{\itshape}% Body font
	{}	% Indent amount
	{\bfseries}	% Theorem head font
	{}	% Punctuation after theorem head
	{.5em}	% Space after theorem head
	{\thmname{#1}\thmnumber{ #2}\thmnote{ \text{(#3)}}}	% head content

\ifcsname c@english\endcsname
	\theoremstyle{plain}
	\newtheorem{theorem}{Theorem}[section]
	\newtheorem{lemma}[theorem]{Lemma}
	\newtheorem{proposition}[theorem]{Proposition}
	\newtheorem*{proposition*}{Proposition}
	%\newtheorem{corollary}[theorem]{אין חלופה עברית}

	\theoremstyle{definition}
	\newtheorem{definition}[theorem]{Definition}
	\newtheorem*{definition*}{Definition}
	\newtheorem{example}{Example}[section]
	\newtheorem{exercise}{Exercise}[section]

	\theoremstyle{remark}
	\newtheorem*{remark}{Remark}
	\newtheorem*{solution}{Solution}
	\newtheorem{conclusion}[theorem]{Conclusion}
	\newtheorem{notation}[theorem]{Notation}
\else
	\theoremstyle{c_plain}
	\newtheorem{theorem}{משפט}[section]
	\newtheorem{lemma}[theorem]{למה}
	\newtheorem{proposition}[theorem]{טענה}
	\newtheorem*{proposition*}{טענה}
	%\newtheorem{corollary}[theorem]{אין חלופה עברית}

	\theoremstyle{c_definition}
	\newtheorem{definition}[theorem]{הגדרה}
	\newtheorem*{definition*}{הגדרה}
	\newtheorem{example}{דוגמה}[section]
	\newtheorem{exercise}{תרגיל}[section]

	\theoremstyle{c_remark}
	\newtheorem*{remark}{הערה}
	\newtheorem*{solution}{פתרון}
	\newtheorem{conclusion}[theorem]{מסקנה}
	\newtheorem{notation}[theorem]{סימון}
\fi

% Questions related commands
\newcounter{question}
\setcounter{question}{1}
\newcounter{sub_question}
\setcounter{sub_question}{1}

\ifcsname c@english\endcsname
	\newcommand{\question}[1][0]{
		\ifthenelse{#1 = 0}{}{\setcounter{question}{#1}}
		\section{Question \arabic{question}}
		\addtocounter{question}{1}
		\setcounter{sub_question}{1}
	}

	\newcommand{\subquestion}[1][0]{
		\ifthenelse{#1 = 0}{}{\setcounter{sub_question}{#1}}
		\subsection{Part \alph{sub_question}}
		\addtocounter{sub_question}{1}
	}
\else
	\newcommand{\question}[1][0]{
		\ifthenelse{#1 = 0}{}{\setcounter{question}{#1}}
		\section{שאלה \arabic{question}}
		\addtocounter{question}{1}
		\setcounter{sub_question}{1}
	}

	\newcommand{\subquestion}[1][0]{
		\ifthenelse{#1 = 0}{}{\setcounter{sub_question}{#1}}
		\subsection{סעיף \localecounter{letters.gershayim}{sub_question}}
		\addtocounter{sub_question}{1}
	}
\fi

% import lua and start of document
\directlua{common = require ('../common')}

\GetEnv{AUTHOR}

% headers
\author{\AUTHOR}
\date\today

\title{פתרון מטלה 09 --- אנליזה פונקציונלית, 80417}

\DeclareMathOperator\sgn{sgn}

\begin{document}
\maketitle
\maketitleprint[teal]

\question{}
תהינה $f, g : \RR \to \RR$ אינטגרביליות רימן ומחזוריות, ונניח ש־$x_0 \in [-\pi, \pi]$.
נניח גם ש־$f \equiv g$ בסביבה של $x_0$.
נסמן $S_N^f, S_N^g$ סכום פורייה מסדר $N$ של $f, g$. \\
נראה ש־$\lim_{N \to \infty} S_N^f(x_0) = L$ אם ורק אם $\lim_{N \to \infty} S_N^g(x_0) = L$.
\begin{proof}
	נניח ש־$\lim_{N \to \infty} S_N^f(x_0) = L$.
	נבחין כי $h = f - g$ מקיימת $h \equiv 0$ בסביבת $x_0$.
	בפרט הגבול שלה ב־$x_0$ הוא $0$ בלבד.
	בנוסף,
	\[
		|h(x_0 + u) - h(x_0)|
		= 0 \le u
	\]
	ונקבל שתנאי ליפשיץ חל, ובהתאם,
	\[
		\lim_{N \to \infty} S_N^{f - g}(x_0)
		= h(x_0)
		= 0
	\]
	אבל אנו יודעים כי $S_N^{f - g} = S_N^f - S_N^g$ מתכונות ולכן למעשה קיבלנו,
	\[
		\lim_{N \to \infty} S_N^f(x_0) - S_N^g(x_0)
		= 0
	\]
	ולכן מההנחה נובע ש־$\lim_{N \to \infty} S_N^g(x_0) = L$ גם כן. \\
	מטעמי סימטריה ההוכחה עבור הכיוון השני זהה.
\end{proof}

\question{}
נוכיח שקיים קבוע $C > 0$ כך שלכל $N > 1$ מתקיים,
\[
	\int_{-\pi}^{\pi} | D_N(u) |\ du
	\ge C \ln(N - 1)
\]
\begin{proof}
	\begin{align*}
		\int_{-\pi}^{\pi} | D_N(u) |\ du
		& = \int_{-\pi}^{\pi} \left\lvert \frac{\sin((N + \frac{1}{2}) u)}{2 \sin \frac{u}{2}} \right\rvert\ du \\
		& = \int_{0}^{\pi} \frac{|\sin((N + \frac{1}{2}) u)|}{2 \sin \frac{u}{2}}\ du - \int_{-\pi}^{0} \frac{|\sin((N + \frac{1}{2}) u)|}{2 \sin \frac{u}{2}}\ du \\
		& = \int_{0}^{\pi} \frac{|\sin((N + \frac{1}{2}) u)|}{\sin \frac{u}{2}}\ du \\
		& \ge \int_{0}^{\pi} \frac{|\sin((N + \frac{1}{2}) u)|}{\frac{u}{2}}\ du \\
		& \ge \int_{0}^{\frac{\pi}{N + \frac{1}{2}}} \frac{1}{\frac{u}{2}}\ du \\
		& \ge 2 \ln(N + 1)
	\end{align*}
\end{proof}

\question{}
\subquestion{}
נראה ש־$A = \Sp {\{ e^{in x} \}}_{n \in \ZZ}$ היא אלגברה ב־$C([-\pi, \pi], \CC)$ שלא מתאפסת באף נקודה ומפרידה נקודות, ואף סגורה להצמדה. \\
נסיק ש־$A$ צפופה ב־$\tilde{C}([-\pi, \pi], \CC)$ בנורמה $\lVert \cdot \rVert_\infty$.
\begin{proof}
	תחילה נראה ש־$A$ אלגברה.
	עבור חיבור הטענה טריוויאלית, עבור כפל,
	\[
		e^{i n x} \cdot e^{i m x}
		= e^{i (n + m) x}
	\]
	והטענה נובעת גם, ועבור סגירות לכפל בסקלר נובע מהגדרת $\Sp$.

	נראה ש־$A$ לא מתאפסת.
	נבחר $f(x) = 1$ על־ידי שימוש ב־$n = 0$.

	נראה ש־$A$ מפרידה נקודות.
	מספיק שנבחר $f(x) = e^{i x}$ ונקבל שמלבד $\pm \pi$ הפונקציה היא חד־חד ערכית ועל.

	מתקיים גם,
	\[
		\overline{e^{i n x}}
		= \cos(n x) - i \sin(n x)
		= \cos(- n x) + i \sin (- n x)
		= e^{-i n x}
	\]
	ולכן $A$ סגורה גם להצמדה.

	כל תנאי משפט סטון־ויירשטראס המרוכב חלים ולכן $A$ צפופה.
\end{proof}

\subquestion{}
נסיק שהקבוצה $A$ היא מערכת אורתוגונלית שלמה ב־$\tilde{C}([-\pi, \pi], \CC)$ יחד עם המכפלה הפנימית,
\[
	\langle f, g \rangle
	= \int_{-\pi}^{\pi} \overline{f}(x) g(x)\ dx 
\]
\begin{proof}
	לכל $n \ne m$,
	\[
		\langle e^{i n x}, e^{i m x} \rangle
		= \int_{-\pi}^{\pi} e^{-i n x} e^{i m x}\ dx
		= \int_{-\pi}^{\pi} e^{i (m - n) x}
		= 0
	\]
	ממשפט קושי למסילות סגורות ורציפות במישור המרוכב נסיק ישירות את הערך הסופי $0$.
	אם כך הסדרה היוצרת את $A$ היא מערכת אורתוגונלית ומאופטימליות טורי פורייה ותנאים שקולים למערכות שלמות נקבל ש־$A$ מערכת שלמה.
\end{proof}

\subquestion{}
עבור $f \in C([-\pi, \pi], \CC)$ נגדיר לכל $n \in \ZZ$ את מקדם פורייה ביחס למערכת המגדירה את $A$ על־ידי,
\[
	c_n
	= \frac{1}{2 \pi} \int_{-\pi}^{\pi} f(x) e^{-i n x}\ dx
\]
נסיק שאם $f$ ממשית אז מתקיים,
\[
	\frac{a_0}{2} = c_0,
	\quad
	\forall n > 0,\ a_n = c_n + c_{-n},
	\quad
	\forall n > 0,\ b_n = i(c_n - c_{-n})
\]
\begin{proof}
	ישירות מהגדרה נקבל,
	\begin{align*}
		a_n
		& = \frac{1}{2 \pi} \int_{-\pi}^{\pi} 2 f(x) \cos(nx)\ dx \\
		& = \frac{1}{2 \pi} \int_{-\pi}^{\pi} f(x) \cos(nx) + i f(x) \sin(nx) - i f(x) \sin(nx) + f(x) \cos(nx)\ dx \\
		& = \frac{1}{2 \pi} \int_{-\pi}^{\pi} f(x) \cos(nx) + i f(x) \sin(nx) + i f(x) \sin(-nx) + f(x) \cos(-nx)\ dx \\
		& = c_n + c_{-n}
	\end{align*}
	ישירות מהגדרה.
	נבחין כי נובע ישירות גם $\frac{a_0}{2} = c_0$ משוויון זה.
	עבור המקרה של $b_n$,
	\begin{align*}
		b_n
		& = \frac{1}{2 \pi} \int_{-\pi}^{\pi} 2 f(x) \sin(nx)\ dx \\
		& = \frac{1}{2 \pi} i \int_{-\pi}^{\pi} i 2 f(x) \sin(-nx) + f(x) \cos(nx) - f(x) \cos(nx)\ dx \\
		& = \frac{1}{2 \pi} i \int_{-\pi}^{\pi} i f(x) \sin(nx) + f(x) \cos(nx) - i f(x) \sin(-nx) - f(x) \cos(-nx)\ dx \\
		& = i(c_n - c_{-n})
	\end{align*}
	וקיבלנו כי השוויון נכון.
\end{proof}

\end{document}
